\documentclass{scrartcl}
\usepackage[utf8]{inputenc}
\usepackage{hyperref}
\usepackage{url}
\usepackage{natbib}
\usepackage{graphicx}
\usepackage{listings}
\usepackage{xcolor}
\usepackage{amsmath}
\usepackage{cleveref}

\lstset{
  basicstyle=\ttfamily\small,
  breaklines=true,
  showstringspaces=false,
  commentstyle=\color{gray},
  keywordstyle=\color{blue},
  stringstyle=\color{red},
  backgroundcolor=\color{white}
}

\newcommand{\emailaddr}[1]{\href{mailto:#1}{\texttt{#1}}}

\title{\LARGE Mancala Game}
\subtitle{Final Report for the Distributed Systems Course 2024/2025}

\author{
  Kasra Sabertehrani \\ \emailaddr{kasra.sabertehrani@studio.unibo.it} \\
    Parisa Sokuti \\ \emailaddr{parisa.sokuti@studio.unibo.it}
}

\date{\today}

\begin{document}

\maketitle

\begin{abstract}
This report presents the design and implementation of a distributed Mancala game system built using Spring Boot and WebSockets.
The system enables two players to play the classic Mancala board game over a network in real-time.
The primary focus is on handling distributed challenges such as player disconnection, session management, inactivity detection,
and state consistency across multiple clients.
The implementation demonstrates core distributed systems concepts including asynchronous communication,
fault tolerance, and timeout mechanisms.
The game is implemented with a clean separation between domain logic (game rules), service layer (business logic),
and presentation layer (WebSocket controllers).
Key features include automatic room creation, player joining, real-time game state synchronization,
graceful disconnection handling with reconnection support, and inactivity-based player forfeit mechanisms.
\end{abstract}

\section*{Disclaimer}

During the preparation of this work, the author(s) used GitHub Copilot to assist with code generation and review.
After using this tool, the author(s) reviewed and edited the content as needed and take(s) full responsibility for the content of the final report/artifact.

\section{Concept}\label{concept}

\subsection{Product Type}

This project is a \textbf{web-based, real-time multiplayer game application} combining:
\begin{itemize}
  \item \textbf{Backend}: RESTful API with WebSocket support for real-time communication
  \item \textbf{Frontend}: HTML/CSS/JavaScript web interface for player interaction
  \item \textbf{Game Engine}: Java-based Mancala game logic implementation
\end{itemize}


\subsection{Why is distribution needed?}

Distribution is essential for this project for the following reasons:

\begin{itemize}
  \item \textbf{Geographical separation}: Players must be able to play from different locations across a network
  \item \textbf{Real-time synchronization}: Game state must be consistent across multiple clients;
    any move by one player must be immediately visible to the other
  \item \textbf{Fault tolerance requirement}: Network failures are inevitable;
    the system must handle player disconnections gracefully and allow reconnection without losing game state
  \item \textbf{Concurrent interactions}: Multiple games must be able to run simultaneously on the same server,
    requiring proper concurrency control
  \item \textbf{Scalability consideration}: Although not a focus of this project,
    the distributed design allows multiple game instances to be load-balanced across servers
\end{itemize}


\section{Requirements Elicitation and Analysis}\label{requirements}

\subsection{Functional Requirements}

\begin{enumerate}
  \item \textbf{FR1: Game Room Management}
  \begin{itemize}
    \item The system must allow a player to create a new game room
    \item \textbf{Acceptance Criterion}: A player can create a room, receive a unique room ID,
      and join as Player 1
  \end{itemize}

  \item \textbf{FR2: Player Joining}
  \begin{itemize}
    \item A second player must be able to join an existing room using its ID
    \item \textbf{Acceptance Criterion}: Player 2 can join a room with Player 1,
      and the game automatically starts
  \end{itemize}

  \item \textbf{FR3: Game Rule Enforcement}
  \begin{itemize}
    \item The system must enforce all Mancala game rules including stone sowing, capturing,
      and turn management
    \item \textbf{Acceptance Criterion}: Moves are validated; invalid moves are rejected;
      game logic is correctly executed
  \end{itemize}

  \item \textbf{FR4: Turn-based Gameplay}
  \begin{itemize}
    \item The system must enforce strict turn-based gameplay where only the current player
      can make a move
    \item \textbf{Acceptance Criterion}: Non-active players cannot make moves;
      turns automatically switch after valid moves
  \end{itemize}

  \item \textbf{FR5: Winner Determination}
  \begin{itemize}
    \item The system must determine the winner when the game ends naturally or by forfeit
    \item \textbf{Acceptance Criterion}: Winner is correctly calculated based on scores or disconnect status
  \end{itemize}

  \item \textbf{FR6: Real-time State Synchronization}
  \begin{itemize}
    \item All game state changes must be immediately broadcast to both players
    \item \textbf{Acceptance Criterion}: Both players see the same game state at all times;
      updates arrive within 100ms
  \end{itemize}

  \item \textbf{FR7: Disconnection Handling}
  \begin{itemize}
    \item When a player disconnects, the game should pause and allow reconnection within a grace period
    \item \textbf{Acceptance Criterion}: Disconnected player can rejoin;
      game resumes from disconnection point
  \end{itemize}

  \item \textbf{FR8: Inactivity Detection}
  \begin{itemize}
    \item If a player is inactive for a specified timeout, the opponent automatically wins
    \item \textbf{Acceptance Criterion}: Inactive player is declared forfeit after timeout
  \end{itemize}

  \item \textbf{FR9: Room Cleanup}
  \begin{itemize}
    \item When a game ends (naturally or by timeout), the room and its resources must be cleaned up
    \item \textbf{Acceptance Criterion}: Room is removed from active rooms;
      no resource leaks occur
  \end{itemize}

  \item \textbf{FR10: Error Handling}
  \begin{itemize}
    \item Invalid operations must return meaningful error messages to the client
    \item \textbf{Acceptance Criterion}: Clients receive error messages describing why operations failed
  \end{itemize}
\end{enumerate}

\subsection{Non-functional Requirements}

\begin{enumerate}
  \item \textbf{NFR1: Consistency}
  \begin{itemize}
    \item Game state must be consistent across all clients at all times
    \item Target: Eventual consistency achieved within 100ms
  \end{itemize}

  \item \textbf{NFR2: Availability}
  \begin{itemize}
    \item The system should remain available even if one player temporarily disconnects
    \item Target: 99\% availability during active games (excluding network failures)
  \end{itemize}


  \item \textbf{NFR3: Fault Tolerance}
  \begin{itemize}
    \item The system must gracefully handle network failures and player disconnections
    \item Target: Detect disconnection within 30 seconds; allow reconnection within 5 minutes
  \end{itemize}

\end{enumerate}

\subsection{Implementation Requirements}

\begin{enumerate}
  \item \textbf{IR1: Programming Language}: Java 17
  \begin{itemize}
    \item \textbf{Justification}: Requirement of Spring Boot framework;
      strong typing and mature ecosystem for distributed systems
  \end{itemize}

  \item \textbf{IR2: Backend Framework}: Spring Boot 4.0.3
  \begin{itemize}
    \item \textbf{Justification}: Provides built-in WebSocket support,
      dependency injection, and scheduling capabilities
  \end{itemize}

  \item \textbf{IR3: Real-time Communication}: WebSocket (STOMP protocol)
  \begin{itemize}
    \item \textbf{Justification}: Enables low-latency, bidirectional communication;
      essential for real-time game synchronization
  \end{itemize}

  \item \textbf{IR4: Frontend}: HTML5, CSS3, JavaScript (Vanilla)
  \begin{itemize}
    \item \textbf{Justification}: No external framework required for this simple use case;
      reduces dependencies
  \end{itemize}

  \item \textbf{IR5: Build Tool}: Maven
  \begin{itemize}
    \item \textbf{Justification}: Provided with Spring Boot; standard for Java projects
  \end{itemize}
\end{enumerate}

\subsection{Glossary of Domain Terms}

\begin{itemize}
  \item \textbf{Pit}: A small cup on the Mancala board containing stones
  \item \textbf{Store}: A larger pit at the end of each player's side, used to collect captured stones
  \item \textbf{Sowing}: The action of distributing stones from a selected pit around the board
  \item \textbf{Capture}: When a stone lands in an empty pit on a player's side,
    allowing the opponent's stones in the opposite pit to be captured
  \item \textbf{Turn}: One player's opportunity to select and sow from a pit
  \item \textbf{Game Room}: A virtual space containing exactly two players and one active game
  \item \textbf{Room ID}: A unique identifier for a game room (e.g., ``1'', ``2'')
  \item \textbf{Player ID}: A unique identifier for a player (UUID generated server-side)
  \item \textbf{WebSocket Session}: A persistent bidirectional connection between client and server
  \item \textbf{Disconnection}: Loss of WebSocket connection due to network failure or client action
  \item \textbf{Reconnection}: Re-establishing a WebSocket connection after disconnection
  \item \textbf{Forfeit}: Automatic win granted to opponent due to inactivity or disconnection timeout
\end{itemize}

\subsection{Relevant Distributed System Features}\label{ds-features}

\subsubsection{Transparency}

\textbf{Relevance}: \textbf{PARTIAL}

The system does NOT attempt to hide distribution from the user:
\begin{itemize}
  \item Users are aware they are playing over a network
  \item If a player disconnects, the system explicitly notifies both players
  \item Reconnection is a visible, user-initiated action
\end{itemize}

However, location transparency is achieved: users don't need to know the physical location of the server.

\subsubsection{Fault Tolerance and Dependability}

\textbf{Relevance}: \textbf{CRITICAL}

This is a primary focus of the system:

\begin{itemize}
  \item \textbf{Availability}: The system remains available even if one player temporarily disconnects
  \item \textbf{Reliability}: Games do not lose state due to network failures
  \item \textbf{Integrity}: Game state cannot be corrupted by out-of-order messages or duplicate messages
  \item \textbf{Graceful degradation}: System continues to function with one player disconnected,
    waiting for reconnection
\end{itemize}

\textbf{Implementation mechanisms}:
\begin{itemize}
  \item Timeout and recovery behavior are summarized in the Design section and detailed in Fault-Tolerance
  \item Synchronization and shared-state control are detailed in Design (Data and Consistency Issues)
  \item Concrete implementation examples are provided in Implementation (Timeout Mechanisms)
\end{itemize}



\section{Design}\label{design}

This section describes how the requirements are met through architectural and design choices.
The design is technology-agnostic and focuses on strategies rather than implementation details.

\subsection{Architecture}\label{architecture}

\subsubsection{Architectural Style}

The system follows a \textbf{Client-Server Architecture with Event-Driven Communication}:

\begin{itemize}
  \item \textbf{Server-side}: Single authoritative game server maintains game state and enforces rules
  \item \textbf{Client-side}: Thin clients (browsers) for input and display only
  \item \textbf{Communication}: Asynchronous event-based (WebSocket) for real-time updates
\end{itemize}


\subsubsection{Conceptual Model}

The system operates with the following main entities and their relationships:

\begin{lstlisting}[language=java]
GameRoom {
  roomId: String
  game: Game
}

Game {
  player1: Player
  player2: Player (nullable)
  board: Board
  gameStatus: GameStatus
  disconnectedPlayerId: String (nullable)
}

Board {
  pits: Pit[]  // 14 pits total
}

Player {
  id: String  // UUID
  name: String
}

Pit {
  stones: int
}
\end{lstlisting}

The game state transitions through the following status sequence:

\begin{lstlisting}
WAITING_FOR_PLAYER_2
  -> (Player 2 joins)
  -> PLAYER_1_TURN
  <-> PLAYER_2_TURN
  -> (or due to disconnect/inactivity)
  -> PAUSED_FOR_RECONNECT
  -> (reconnect success)
  -> PLAYER_1_TURN or PLAYER_2_TURN
  -> (or timeout expires)
  -> PLAYER_DISCONNECTED
  -> (game ends)
  -> GAME_OVER
\end{lstlisting}

\subsection{Infrastructure}\label{infrastructure}

\subsubsection{Components}

The system comprises the following infrastructure components:

\begin{enumerate}
  \item \textbf{Web Server (Spring Boot)}:
  \begin{itemize}
    \item Handles HTTP requests for room creation and joining
    \item Exposes WebSocket endpoint for game communication
    \item Single instance per deployment
  \end{itemize}

  \item \textbf{WebSocket Broker}:
  \begin{itemize}
    \item Spring's simple message broker
    \item Routes messages to subscribed clients
    \item Built-in to Spring Boot; no external message queue required
  \end{itemize}

  \item \textbf{Web Browsers (Clients)}:
  \begin{itemize}
    \item Establish WebSocket connections to the server
    \item Send game moves and receive state updates
    \item Multiple browsers per game (exactly 2 players)
  \end{itemize}

  \item \textbf{In-Memory Store}:
  \begin{itemize}
    \item ConcurrentHashMap storing active game rooms
    \item ConcurrentHashMap storing session-to-player mappings
    \item No persistent database (games are ephemeral)
  \end{itemize}
\end{enumerate}



\subsection{Modelling}\label{modelling}

\subsubsection{Domain Entities}

The system models four primary domain entities:

\begin{enumerate}
  \item \textbf{Game}: Represents a single game instance
  \begin{itemize}
    \item Encapsulates game rules and state transitions
    \item Manages both players and the board
    \item Determines winners and enforces turn order
  \end{itemize}

  \item \textbf{Board}: Represents the Mancala board
  \begin{itemize}
    \item 14 pits (6 per player + 1 store each)
    \item Handles stone sowing, capturing, and scoring
    \item Completely decoupled from players or game status
  \end{itemize}

  \item \textbf{Player}: Represents a player in the game
  \begin{itemize}
    \item Server-assigned UUID for identification
    \item Player name provided by the user
    \item No mutable state (purely identity)
  \end{itemize}

  \item \textbf{GameRoom}: Container for a game instance
  \begin{itemize}
    \item Groups a game with its room ID
    \item Minimal business logic; mostly a data holder
  \end{itemize}
\end{enumerate}

\subsubsection{Mapping Domain Entities to Infrastructure}

\textbf{Where game state lives}:

\begin{itemize}
  \item \textbf{Game state (Game, Board, Pits)}: Lives exclusively on the server in memory
  \item \textbf{Player input (moves)}: Sent from client to server via WebSocket
  \item \textbf{Game display state}: Lives on client in JavaScript variables (mirrors server)
  \item \textbf{Session information}: Stored in server's SessionTracker
  \item \textbf{Timing information}: Stored in server's activity/disconnection timestamp maps
\end{itemize}

\textbf{Synchronization strategy}:

The server is the single source of truth and clients act as renderers of server state.
Detailed consistency guarantees are defined in \textbf{Consistency Model} to avoid repeating them here.

\subsubsection{Domain Events}

The system recognizes the following domain events:

\begin{enumerate}
  \item \textbf{RoomCreated}: Player 1 creates a new game room
  \item \textbf{PlayerJoined}: Player 2 joins an existing room, game starts
  \item \textbf{MoveMade}: A player selects a pit and makes a valid move
  \item \textbf{TurnSwitched}: The turn automatically changes to the other player
  \item \textbf{PlayerDisconnected}: A player's WebSocket connection is lost
  \item \textbf{PlayerReconnected}: A disconnected player re-establishes connection
  \item \textbf{ReconnectionTimeoutExpired}: Grace period for reconnection expires; opponent wins
  \item \textbf{PlayerInactive}: A player makes no moves for the inactivity timeout duration
  \item \textbf{GameEnded}: Game ends naturally (all pits on one side empty)
  \item \textbf{RoomCleaned}: Room is removed from active rooms
\end{enumerate}

\subsubsection{Message Types}

The system exchanges the following types of messages over WebSocket:

\begin{enumerate}
  \item \textbf{Commands} (Client -> Server):
  \begin{itemize}
    \item \texttt{/app/game.move}: Execute a move (pit index, room ID, player ID)
    \item \texttt{/app/game.join}: Join a game (room ID, player ID)
    \item \texttt{/app/game.reconnect}: Attempt to reconnect (room ID, player ID)
  \end{itemize}

  \item \textbf{Events} (Server -> Client, broadcast):
  \begin{itemize}
    \item \texttt{Game state update}: Broadcasted to \texttt{/topic/room/{roomId}}
    \item Contains: Player names, pit states, current turn, game status, scores, winner, disconnection info
  \end{itemize}

  \item \textbf{Errors} (Server -> Client, point-to-point):
  \begin{itemize}
    \item \texttt{/queue/errors}: Invalid moves, room not found, invalid state errors
  \end{itemize}
\end{enumerate}



\subsection{Interaction}\label{interaction}

\subsubsection{Communication Patterns}

The system employs two main communication patterns:

\textbf{1. Request-Response (HTTP)}

Used for game room creation and joining:

\begin{lstlisting}
Client                          Server
  |                              |
  +-> POST /api/rooms/create --->|
  |                              |
  |  (Create Player 1)           |
  |  (Create GameRoom)           |
  |  (Return room details)       |
  |<--------- 200 OK ------------+
  |  (room ID, player ID)        |
  |                              |
\end{lstlisting}

\textbf{2. Event-Driven (WebSocket)}

Used for game moves and state synchronization:

\begin{lstlisting}
Player 1                  Server              Player 2
  |                       |                   |
  +-> CONNECT /mancala-ws>|                   |
  |<- CONNECTED ----------+                   |
  |                       |                   |
  +-> /app/game.move ---->| (Validate move)   |
  |                       | (Update state)    |
  |                       +-> /topic/room/1 ->| (Broadcast)
  |<- Game state -----------------------------+|
  |                       |                   |
  | (Player 2's turn)     |                   |
  |                       |<-- /app/game.move +
  |                       | (Validate move)   |
  |<- /topic/room/1 -------------------------+| (Broadcast)
  |   (Game state)        |                   |
  |                       |                   |
\end{lstlisting}

\subsubsection{Timeout Mechanism}

\begin{lstlisting}
Event: Player Disconnects
|
+- Detect disconnect and pause game
+- Start reconnection grace timer
+- If player reconnects in time -> resume
+- Otherwise -> forfeit, cleanup, broadcast final state
\end{lstlisting}

Detailed timeout values and scheduler behavior are specified in Fault-Tolerance and Implementation.

\subsection{Behaviour}\label{behaviour}

\subsubsection{Game State Machine}

The game transitions through discrete states as follows:

\textbf{State: WAITING\_FOR\_PLAYER\_2}
\begin{itemize}
  \item \textbf{Entry}: Player 1 creates a room
  \item \textbf{Actions}: Server waits for Player 2 to join
  \item \textbf{Transitions}:
  \begin{itemize}
    \item -> PLAYER\_1\_TURN when Player 2 joins
  \end{itemize}
  \item \textbf{Constraints}: Only Player 1 exists; no moves allowed
\end{itemize}

\textbf{State: PLAYER\_1\_TURN}
\begin{itemize}
  \item \textbf{Entry}: Game starts or after Player 2's move
  \item \textbf{Actions}: Player 1 can select a pit and move
  \item \textbf{Transitions}:
  \begin{itemize}
    \item -> PLAYER\_1\_TURN if Player 1 lands in their store (extra turn)
    \item -> PLAYER\_2\_TURN if Player 1's turn ends normally
    \item -> PAUSED\_FOR\_RECONNECT if Player 1 disconnects
    \item -> GAME\_OVER if game ends
  \end{itemize}
\end{itemize}

\textbf{State: PLAYER\_2\_TURN}
\begin{itemize}
  \item \textbf{Entry}: After Player 1's turn ends (normal transition)
  \item \textbf{Actions}: Player 2 can select a pit and move
  \item \textbf{Transitions}: (Symmetric to PLAYER\_1\_TURN)
\end{itemize}

\textbf{State: PAUSED\_FOR\_RECONNECT}
\begin{itemize}
  \item \textbf{Entry}: A player's WebSocket disconnects during an active turn
  \item \textbf{Actions}: Server remembers the previous state and waits for reconnection
  \item \textbf{Transitions}:
  \begin{itemize}
    \item -> PLAYER\_1\_TURN or PLAYER\_2\_TURN if disconnected player reconnects (resume)
    \item -> PLAYER\_DISCONNECTED if reconnection timeout expires
  \end{itemize}
  \item \textbf{Constraints}: No new moves allowed; waiting player cannot act
\end{itemize}

\textbf{State: PLAYER\_DISCONNECTED}
\begin{itemize}
  \item \textbf{Entry}: Reconnection timeout expires without player reconnecting
  \item \textbf{Actions}: Game is forfeit; opponent automatically wins
  \item \textbf{Transitions}:
  \begin{itemize}
    \item -> GAME\_OVER (final state)
  \end{itemize}
\end{itemize}

\textbf{State: GAME\_OVER}
\begin{itemize}
  \item \textbf{Entry}: Game ends naturally (board empty on one side) or by forfeit
  \item \textbf{Actions}: Calculate winner; return game to clients
  \item \textbf{Transitions}: (Terminal state)
  \item \textbf{Cleanup}: Room is removed after broadcast
\end{itemize}


\subsection{Data and Consistency Issues}\label{data-and-consistency-issues}

\subsubsection{Data Storage}

\textbf{What data is stored?}

\begin{itemize}
  \item Game state: Pit values, turn information, game status
  \item Player information: Names, UUIDs
  \item Room information: Room IDs, room -> game mappings
  \item Session information: WebSocket session IDs -> player/room mappings
  \item Temporal information: Activity timestamps, disconnection timestamps
\end{itemize}

\textbf{Where is it stored?}

\begin{itemize}
  \item All data in server's Java heap memory
  \item Two primary structures: \texttt{activeRooms} (ConcurrentHashMap) and \texttt{sessionSessions} (ConcurrentHashMap)
  \item No persistent storage (no database)
\end{itemize}


\subsubsection{Consistency Model}

The system implements \textbf{Strong Consistency} for game state:

\begin{itemize}
  \item At any point in time, there is exactly one version of the game state
  \item All moves are applied sequentially (never in parallel)
  \item Before returning a move result, the new state is immediately broadcast to all clients
  \item Clients do not make independent decisions; they always display server-truth
\end{itemize}

\textbf{Consistency guarantees}:

\begin{itemize}
  \item No lost updates: Each move is persisted in server memory immediately
  \item No dirty reads: Clients only read state after server confirms
  \item No race conditions: Synchronization prevents concurrent state modifications
  \item Ordered updates: Messages are delivered in the order they were sent
\end{itemize}

\subsubsection{Query Operations}

\textbf{Who queries the database?}

Since there is no persistent database, we consider in-memory queries:

\begin{itemize}
  \item \textbf{RoomService}: Queries activeRooms map when retrieving/joining/removing rooms
  \item \textbf{GameService}: Queries rooms when processing moves and timeouts
  \item \textbf{SessionTracker}: Queries session mappings when tracking/removing sessions
\end{itemize}

\textbf{What queries?}

\begin{itemize}
  \item Get room by ID: O(1) hash lookup
  \item Get all active rooms: O(n) iteration (rarely needed)
  \item Check session membership: O(1) hash lookup
\end{itemize}

\textbf{When?}

\begin{itemize}
  \item On every move: ~1ms processing time
  \item On disconnection event: Immediate
  \item On timeout check: Every 5 seconds (scheduled)
\end{itemize}

\subsubsection{Shared State Management}

The following data structures are shared between multiple components:

\begin{itemize}
  \item \textbf{GameRoom}: Shared between RoomService (creator), GameService (updater), WebSocketEventListener (accessor)
  \item \textbf{Session mappings}: Shared between GameController (writer) and WebSocketEventListener (reader)
  \item \textbf{Activity timestamps}: Shared between GameService (writer) and GameTimeoutScheduler (reader)
\end{itemize}

Thread safety is ensured through:

\begin{itemize}
  \item ConcurrentHashMap for all shared mutable collections
  \item Synchronized blocks around complex multi-step operations
  \item Immutable value objects (Player, Pit are largely immutable)
\end{itemize}

\subsection{Fault-Tolerance}\label{fault-tolerance}


\subsubsection{Heartbeat and Timeout Mechanisms}

\textbf{WebSocket-level heartbeat}:

\begin{itemize}
  \item SockJS (WebSocket fallback) provides automatic heartbeating
  \item Server detects client disconnection within a few seconds
  \item Triggered by network issues or client going offline
\end{itemize}

\textbf{Application-level timeouts}:

\begin{enumerate}
  \item \textbf{Reconnection grace period}: 30 seconds (resume if player returns; otherwise forfeit)
  \item \textbf{Inactivity timeout}: 60 seconds (current player forfeits when idle beyond threshold)
  \item \textbf{Scheduler check interval}: 5 seconds (periodically evaluates both timeout rules)
\end{enumerate}

\subsubsection{Error Handling}

\textbf{Move validation errors}:

\begin{lstlisting}
Invalid pit index --> IllegalArgumentException
  |
GameController.@MessageExceptionHandler
  |
Send error to client via /queue/errors
  |
Client displays error message
\end{lstlisting}

\textbf{Game state errors}:

\begin{itemize}
  \item Move on wrong turn: Throw IllegalStateException
  \item Move on game over: Throw IllegalStateException
  \item Invalid player ID: Throw IllegalArgumentException
\end{itemize}

\textbf{Room not found errors}:

\begin{itemize}
  \item Room doesn't exist: Return HTTP 404
  \item Room is full: Return HTTP 400
  \item Proper error messages sent to client
\end{itemize}

\subsubsection{Recovery Strategy}

\textbf{From temporary disconnection (< 30 seconds)}:

\begin{itemize}
  \item \textbf{Server-Side:} The server pauses the game state and waits for a STOMP reconnect payload containing the original Player ID.
  \item \textbf{Client-Side Recovery:} The frontend utilizes both \texttt{localStorage} (to save the room and player IDs) and \texttt{sessionStorage} (to verify if the player was actively in a match within that specific tab). If a user accidentally refreshes the page, the JavaScript \texttt{window.onload} event intercepts the load, reads the \texttt{localStorage} payload, bypasses the lobby screen, and automatically fires an \texttt{/app/game.reconnect} message to the server to instantly resume the session.
\end{itemize}

\textbf{From permanent disconnection (> 30 seconds)}:

\begin{itemize}
  \item Opponent is declared winner
  \item Room is cleaned up
  \item Other player is notified
  \item Disconnected player sees game over state if they reconnect late
\end{itemize}

\subsection{Availability}\label{availability}

\subsubsection{Caching}

\textbf{Client-side caching}:

\begin{itemize}
  \item Player ID and room ID stored in browser localStorage
  \item Used to attempt reconnection to the same game
  \item Reduces need to re-enter information after page refresh
\end{itemize}

\textbf{Server-side caching}:

\begin{itemize}
  \item No caching; all state is in-memory anyway
  \item No database to cache query results
\end{itemize}

\subsubsection{Load Balancing}

\textbf{Not implemented} for this project:

\begin{itemize}
  \item Single server instance handles all games
  \item If many games run simultaneously, server CPU/memory could saturate
  \item Production version would use load balancer to distribute across multiple servers
\end{itemize}

\subsubsection{Network Partitioning}

\textbf{Scenario 1: Client-Server partition (client isolated)}

\begin{itemize}
  \item Client detects disconnection
  \item Game pauses on server (PAUSED\_FOR\_RECONNECT)
  \item Opponent continues to see paused game
  \item If partition persists > 30 seconds, game forfeits to opponent
  \item Client cannot send moves (isolated)
\end{itemize}

\textbf{Scenario 2: Client-Server partition (server isolated)}

\begin{itemize}
  \item Server detects both clients disconnected
  \item Server cannot send broadcasts (no recipients)
  \item All games timeout and are cleaned up
  \item Inbound moves are dropped (no clients)
\end{itemize}

\textbf{Scenario 3: Slow network (high latency)}

\begin{itemize}
  \item Moves still processed correctly (all happens on server)
  \item Broadcasts are delayed reaching clients
  \item Game state will eventually converge (eventual consistency achieved)
\end{itemize}

The system favors \textbf{strong consistency} over availability: if there's any doubt about state,
the server pauses the game rather than allowing conflicting moves.

\subsection{Security}\label{security}


\subsubsection{Input Validation}

\textbf{Server-side validation} (strong):

\begin{itemize}
  \item Player ID is validated on every move
  \item Pit index is validated (must be 0-5 for P1, 7-12 for P2)
  \item Move is validated against game rules
  \item Game state is validated (only allow moves in playable states)
\end{itemize}

\textbf{Client-side validation} (weak):

\begin{itemize}
  \item The Vanilla JS actively disables click events and applies CSS \texttt{.disabled} and \texttt{.opponent-pit} classes to prevent players from clicking empty pits or opponent pits.
  \item While not a strict security measure (as it can be bypassed via the browser console), this prevents accidental malformed payloads from reaching the server and provides immediate user experience constraints.
\end{itemize}



\section{Implementation}\label{implementation}

This section describes the technology-dependent choices made to realize the design described above.

\subsection{Network Protocols}

\textbf{Protocol Choice: STOMP over WebSocket (with SockJS fallback)}

\begin{itemize}
  \item \textbf{HTTP REST}: Session setup (room create/join) via \texttt{/api/rooms/create} and \texttt{/api/rooms/join}
  \item \textbf{WebSocket + STOMP}: Real-time gameplay events and state broadcasts with low latency
  \item \textbf{SockJS}: Transport fallback for environments where native WebSocket is unavailable
\end{itemize}


\subsection{Data Representation}

\textbf{In-transit data format: JSON}

\begin{lstlisting}
Example game state broadcast:

{
  "player1": {
    "id": "550e8400-e29b-41d4-a716-446655440000",
    "name": "Alice"
  },
  "player2": {
    "id": "6ba7b810-9dad-11d1-80b4-00c04fd430c8",
    "name": "Bob"
  },
  "board": {
    "pits": [4, 4, 4, 4, 4, 4, 0, 4, 4, 4, 4, 4, 4, 0]
  },
  "gameStatus": "PLAYER_1_TURN",
  "player1Score": 0,
  "player2Score": 0,
  "disconnectedPlayerId": null,
  "winner": null
}
\end{lstlisting}


\subsection{In-Memory Data Structures}

\begin{lstlisting}[language=java]
// Room storage
ConcurrentHashMap<String, GameRoom> activeRooms

// Session tracking
ConcurrentHashMap<String, SessionInfo> sessions
  where SessionInfo = (playerId, roomId)

// Activity tracking
ConcurrentHashMap<String, Instant> lastActivityTimestamps
ConcurrentHashMap<String, Instant> disconnectTimestamps

// Synchronization
synchronized (room) { ... }
\end{lstlisting}

\textbf{Why ConcurrentHashMap?}

\begin{itemize}
  \item Thread-safe without locking entire map
  \item Multiple readers can access simultaneously
  \item Writers use fine-grained locks (lock striping)
  \item Perfect for multiple concurrent games
\end{itemize}

\textbf{Why synchronized blocks?}

\begin{itemize}
  \item Ensure complex multi-step operations are atomic
  \item Prevent inconsistent state during game transitions
  \item Example: Move validation + board update + turn switch must be indivisible
\end{itemize}

\subsection{Timeout Mechanisms}

This section contains implementation-level examples only. Policy definitions (values and rationale)
are provided earlier in Design and Fault-Tolerance.

\textbf{Reconnection timeout implementation example}

\begin{lstlisting}[language=java]
private static final Duration RECONNECT_GRACE = Duration.ofSeconds(30);

// When player disconnects:
disconnectTimestamps.put(roomId, Instant.now());

// In scheduler (every 5 seconds):
Duration elapsed = Duration.between(disconnectedAt, Instant.now());
if (elapsed.compareTo(RECONNECT_GRACE) > 0) {
  // Forfeit game
  room.getGame().forfeit(disconnectedPlayerId);
}
\end{lstlisting}

\textbf{Inactivity timeout implementation example}

\begin{lstlisting}[language=java]
private static final Duration INACTIVITY_TIMEOUT = Duration.ofMinutes(1);

// When move is made:
lastActivityTimestamps.put(roomId, Instant.now());

// In scheduler (every 5 seconds):
Duration inactive = Duration.between(lastMove, Instant.now());
if (inactive.compareTo(INACTIVITY_TIMEOUT) > 0) {
  // Current player is idle; they lose
  String idlePlayer = (gameStatus == PLAYER_1_TURN)
    ? player1.getId()
    : player2.getId();
  room.getGame().forfeit(idlePlayer);
}
\end{lstlisting}

\textbf{Scheduler invocation implementation example}

\begin{lstlisting}[language=java]
@Scheduled(fixedRate = 5000)  // 5000 milliseconds
public void checkTimeouts() {
  List<GameRoom> timedOutRooms = gameService.processTimeouts();
  for (GameRoom room : timedOutRooms) {
    messagingTemplate.convertAndSend("/topic/room/" + room.getRoomId(), room.getGame());
  }
}
\end{lstlisting}

\subsection{Technological Details}\label{technological-details}

\subsubsection{Spring Boot Framework}

\textbf{Version}: 4.0.3

\textbf{Key dependencies}:

\begin{itemize}
  \item \textbf{spring-boot-starter-webmvc}: Provides HTTP server and REST controllers
  \item \textbf{spring-boot-starter-websocket}: Enables WebSocket + STOMP support
  \item \textbf{spring-boot-starter-webmvc-test}: Testing framework
  \item \textbf{lombok}: Reduces boilerplate (not heavily used in this project)
\end{itemize}

\subsection{Frontend Implementation (Client-Side)}

The client-side interface was developed exclusively using HTML5, CSS3, and Vanilla JavaScript, completely avoiding heavy frameworks to ensure a lightweight, instantly loading application.

\textbf{User Interface and Styling:}
\begin{itemize}
  \item \textbf{DOM Structure:} The HTML is divided into two primary views: the \texttt{lobby-screen} for room creation and the \texttt{game-screen} for the active board. The transition between these screens is handled dynamically via CSS classes.
  \item \textbf{CSS Animations:} Visual cues are strictly enforced using CSS. For example, when it is a player's turn, their valid pits receive a \texttt{.active} class that triggers a custom \texttt{@keyframes pulse} animation, drawing the user's attention. Opponent pits are visually locked with a \texttt{.disabled} class and a restricted cursor.
\end{itemize}

\textbf{WebSocket Integration (app.js):}
\begin{itemize}
  \item \textbf{STOMP \& SockJS:} The client utilizes the SockJS library to establish the connection and the STOMP protocol to format the messages.
  \item \textbf{Event Listeners:} The JavaScript attaches click event listeners to all \texttt{.pit} elements. When clicked, it extracts the pit index (e.g., parsing ``4'' from \texttt{pit-4}), verifies if it is locally the player's turn to prevent spam clicking, and sends a JSON payload to the server's \texttt{/app/game.move} endpoint.
  \item \textbf{State Rendering:} Upon receiving a broadcast from \texttt{/topic/room/\{roomId\}}, the \texttt{updateBoard()} function iterates through the JSON state, updating the inner text of all 14 pits, checking the \texttt{gameStatus} to update the UI banners (e.g., switching to a ``Paused'' state), and enabling/disabling pit interactivity.
\end{itemize}


\section{Validation}\label{validation}

\subsubsection{Manual Testing (Performed)}

The following manual testing was performed during development. Visual evidence of these tests can be seen in \cref{fig:manual-in-game,fig:manual-disconnection,fig:manual-disconnection-2,fig:manual-back-in-game,fig:manual-inactivity}.

\begin{enumerate}
  \item \textbf{Normal game flow}
  \begin{itemize}
    \item Created room
    \item Joined room from second browser
    \item Played complete game
    \item Verified winner determination
    \item \textbf{Result}: PASS
  \end{itemize}

  \item \textbf{Disconnection and reconnection}
  \begin{itemize}
    \item Simulated disconnection by blocking network (DevTools)
    \item Waited < 30 seconds
    \item Reconnected
    \item Game resumed correctly
    \item \textbf{Result}: PASS
  \end{itemize}

  \item \textbf{Reconnection timeout}
  \begin{itemize}
    \item Simulated disconnection
    \item Waited > 30 seconds
    \item Opponent was declared winner
    \item \textbf{Result}: PASS
  \end{itemize}

  \item \textbf{Inactivity timeout}
  \begin{itemize}
    \item Set inactivity timeout to 60 seconds (for testing)
    \item Played moves until active player stopped moving
    \item Waited > 60 seconds
    \item Inactive player was forfeit
    \item \textbf{Result}: PASS
  \end{itemize}

\end{enumerate}

\textbf{Manual testing screenshots:}

\begin{figure}
  \centering
  \includegraphics[width=0.85\textwidth]{\detokenize{figures/inside the game.png}}
  \caption{Manual test evidence: normal in-game state after both players join.}
  \label{fig:manual-in-game}
\end{figure}

\begin{figure}
  \centering
  \includegraphics[width=0.85\textwidth]{\detokenize{figures/disconnection.png}}
  \caption{Manual test evidence: disconnection detected and reconnect wait state shown.}
  \label{fig:manual-disconnection}
\end{figure}

\begin{figure}
  \centering
  \includegraphics[width=0.85\textwidth]{\detokenize{figures/disconnection-2.png}}
  \caption{Manual test evidence: reconnection timeout progression.}
  \label{fig:manual-disconnection-2}
\end{figure}

\begin{figure}
  \centering
  \includegraphics[width=0.85\textwidth]{\detokenize{figures/back in game.png}}
  \caption{Manual test evidence: successful recovery after reconnection.}
  \label{fig:manual-back-in-game}
\end{figure}

\begin{figure}
  \centering
  \includegraphics[width=0.85\textwidth]{\detokenize{figures/inactivity.png}}
  \caption{Manual test evidence: inactivity timeout and automatic forfeit outcome.}
  \label{fig:manual-inactivity}
\end{figure}


\section{Deployment}\label{deployment}

\subsection{Installation Instructions}

\subsubsection{Prerequisites}

\begin{itemize}
  \item Java 17 or later installed
  \item Maven 3.6+ installed
  \item Git (to clone repository)
  \item Modern web browser (Chrome, Firefox, Safari, Edge)
\end{itemize}

\subsubsection{Build and Run}

\begin{lstlisting}[language=bash]
# 1. Clone the repository
git clone <repository-url>
cd mancala-game-springboot

# 2. Build the project
mvn clean package

# 3. Run the application
mvn spring-boot:run


# 4. Access in browser
# Navigate to http://localhost:8080
\end{lstlisting}



\section{User Guide}\label{user-guide}

\subsection{Starting the Application}

\begin{enumerate}
  \item Open terminal/command prompt
  \item Navigate to project directory
  \item Run: \texttt{mvn spring-boot:run}
  \item Wait for startup message: ``Started MancalaGameSpringbootApplication in X seconds''
  \item Open browser to \texttt{http://localhost:8080}
\end{enumerate}

\subsection{Creating a Game}

\begin{enumerate}
  \item On the home page, enter your player name
  \item Click ``Create Room''
  \item A room ID will be displayed (e.g., ``1'')
  \item Share this room ID with your opponent (via chat, email, etc.)
  \item Game waits with message ``Waiting for Player 2''
\end{enumerate}

\subsection{Joining a Game}

\begin{enumerate}
  \item Obtain room ID from Player 1
  \item Open browser to \texttt{http://localhost:8080}
  \item Enter your player name
  \item Enter the room ID
  \item Click ``Join Room''
  \item Game board appears; Player 1 makes first move
\end{enumerate}

\subsection{Playing the Game}

\begin{itemize}
  \item Click any pit on your side to select and sow stones
  \item The system enforces all Mancala rules:
  \begin{itemize}
    \item Stones are distributed around the board
    \item Captures occur when landing in empty pits
    \item Extra turns granted when landing in your store
    \item Game ends when one side is empty
  \end{itemize}
  \item Game status is displayed at top (whose turn)
  \item Scores update in real-time
  \item Winner is announced when game ends
\end{itemize}

\subsection{Handling Disconnection}

\begin{itemize}
  \item If your connection is lost, game pauses
  \item Message appears: ``Opponent disconnected''
  \item You have 30 seconds to reconnect before opponent wins automatically
  \item To reconnect: Refresh page; system detects and resumes game
  \item If reconnection fails within 30 seconds, game is forfeited
\end{itemize}

\subsection{Rules Reminder}

\textbf{How Mancala is played}:

\begin{enumerate}
  \item Players alternate turns
  \item On your turn, select a pit on your side
  \item Stones are distributed counter-clockwise
  \item Stones are skipped in opponent's store
  \item If last stone lands in your store, you get another turn
  \item If last stone lands in an empty pit on your side, capture opposite pit
  \item Game ends when one side is completely empty
  \item Winner is player with most stones in their store
\end{enumerate}

\section{Self-evaluation}\label{self-evaluation}

\subsection{Kasra Sabertehrani}

\textbf{Role within the group:}
As the backend developer, my primary responsibility was engineering the authoritative server and managing the distributed state. I designed the Spring Boot architecture, implemented the STOMP WebSocket message broker, and wrote the core Mancala game engine. I specifically focused on solving concurrency challenges by designing thread-safe operations using \texttt{ConcurrentHashMap} and synchronized blocks to ensure state consistency during simultaneous network requests. Additionally, I implemented the \texttt{SessionTracker} and the scheduled timeout mechanisms to handle silent TCP connection drops.

\textbf{Product Strengths (Backend Perspective):}
\begin{itemize}
  \item \textbf{Reliable Real-Time Sync}: The WebSocket broadcasting infrastructure successfully keeps both players aligned with sub-second latency.
  \item \textbf{Safe Concurrency}: The explicit thread locking prevents race conditions, successfully supporting multiple simultaneous game rooms without cross-room interference.
  \item \textbf{Robust Session Handling}: Disconnection/reconnection and timeout paths are explicitly modeled, preventing the server from holding onto dead connections indefinitely.
\end{itemize}

\textbf{Product Weaknesses (Backend Perspective):}
\begin{itemize}
  \item \textbf{Ephemeral State}: No persistent database is utilized; all active games and player sessions are lost if the Spring Boot server restarts.
  \item \textbf{Single Server Deployment}: The architecture lacks horizontal scaling. If the single node crashes, there is no failover mechanism.
  \item \textbf{Security}: Communication is unencrypted (no HTTPS/WSS) and lacks a formal authentication protocol.
\end{itemize}

\subsection{Parisa Sokuti}

\textbf{Role within the group:}
As the frontend developer, my role focused on building the client-side architecture and user interface using HTML, CSS, and Vanilla JavaScript. I was responsible for integrating the SockJS and STOMP client libraries to establish the persistent connection with the server. I engineered the logic to parse incoming JSON state broadcasts and dynamically manipulate the DOM to render the Mancala board, update scores, and provide real-time visual feedback for game events, including network state changes (e.g., displaying "Paused" screens when the opponent disconnects).

\textbf{Product Strengths (Frontend Perspective):}
\begin{itemize}
  \item \textbf{Lightweight Client}: By utilizing Vanilla JS rather than a heavy framework like React or Angular, the client application loads instantly and minimizes unnecessary overhead.
  \item \textbf{Responsive Feedback}: The UI reacts strictly to server-driven events, ensuring the player's view is always perfectly synchronized with the authoritative state.
  \item \textbf{Graceful Degradation}: The interface clearly communicates network issues to the user, masking the complexity of the distributed system behind user-friendly countdowns and pause screens.
\end{itemize}

\textbf{Product Weaknesses (Frontend Perspective):}
\begin{itemize}
  \item \textbf{DOM Manipulation Complexity}: Direct DOM manipulation in Vanilla JS becomes increasingly difficult to maintain as the application's complexity grows compared to a component-based framework.
  \item \textbf{Client-Side Caching}: If a player refreshes the page manually, the UI completely resets until the server pushes the next full state broadcast, causing a momentary flash of unstyled or empty content.
  \item \textbf{Basic Styling}: The interface is functional but lacks advanced graphical polish, animations, or responsive mobile optimizations.
\end{itemize}



\section{Future works}\label{future-works}


\begin{enumerate}
  \item \textbf{Spectator Mode}: Allow additional players to watch games without participating

  \item \textbf{Undo/Replay}: Allow players to undo recent moves or replay games

  \item \textbf{AI Opponent}: Implement an AI player for single-player mode

  \item \textbf{Game Variants}: Support different Mancala variants (different board sizes, rules)

  \item \textbf{Chat Feature}: Allow players to communicate during games

  \item \textbf{Rating System}: Track player ratings, history, statistics

  \item \textbf{Tournament Mode}: Support bracketed tournaments with multiple games
\end{enumerate}


\end{document}

